\documentclass[10pt,a4paper]{amsart}
\usepackage[margin=19mm]{geometry}
\usepackage{amsmath,amssymb,mathtools}
\usepackage[scr=rsfs,cal=euler]{mathalfa}
\usepackage{xfrac}
\usepackage[unicode,hidelinks,bookmarks=false]{hyperref}
\newcommand{\ie}{i.e.\ }
\newcommand{\cf}{cf.\ }
\newcommand{\resp}{respectively\ }
\newcommand{\Subsection}{\subsection}
\providecommand{\nobreakdash}{\nobreak\hbox{-}}
\setlength{\emergencystretch}{2em}
\multlinegap0pt
\usepackage{amssymb}
\usepackage{mathtools}
\usepackage{xcolor}
\newtheorem{lemma}{Lemma}
\newtheorem{proposition}{Proposition}
\newcommand{\E}{\mathbb{E}}
\newcommand{\abs}[1]{\left\lvert #1\right\rvert}
\newcommand{\norm}[1]{\left\lVert #1\right\rVert}
\title{Logarithmic-Free Moment and Generalization Bounds for Uniformly Stable Algorithms}
\author{Thanh Nguyen-Cung, Binh T. Nguyen}
\date{Source: arXiv:2608.09870v2 (CC BY 4.0)}
\begin{document}
\maketitle

\noindent\emph{Source and licence.} Logarithmic-Free Moment and Generalization Bounds for Uniformly Stable Algorithms. Version arXiv:2608.09870v2 (CC BY 4.0), distributed by arXiv under the Creative Commons Attribution 4.0 International licence, http://creativecommons.org/licenses/by/4.0/, as recorded in the arXiv licence metadata for this exact version and shown on the abstract page https://arxiv.org/abs/2608.09870v2. Full licence notice: \texttt{licenses/arxiv-2608.09870-CC-BY-4.0.txt}.\par\smallskip

\noindent\textit{Editorial scope.} Counted unit: the article's proposition prop:transfer (source0.tex lines 440--452). The article's complete original proof of it is delivered with it (source0.tex lines 454--554, one \texttt{proof} environment of 101 lines). No mathematics is written, repaired or re-derived: the statement and the proof are the article's own lines, in the article's own notation, and no retained line is substituted or re-flowed.  Retained uncounted as the source-local context the proof needs: lines 244--255; lines 307--316.  Nothing was omitted from any retained span, so the package carries no omission record.  No retained line contains a citation, so the wrapper carries no bibliography and no bibliography entry.  No retained line contains a figure environment, an image-inclusion command or drawing code, so no image is delivered and the package declares no asset dependency.  Editorial formatting only, in addition: an AMS \texttt{amsart} preamble that restates the article's own declarations and macros used by the retained lines, namely the environments \texttt{lemma}, \texttt{proposition} on separate counters, together with the standard packages those lines need.  The retained environments print in this document as Lemma 1, Proposition 1 in order of appearance, whereas the article prints the same items with the numbering of its own full document.  The complete native source of the article as distributed in the arXiv e-print for this exact version is bundled unchanged as \texttt{sources/207279/source0.tex}.
\begin{lemma}\label{lem:bd}
Let $X=f(W_1,\ldots,W_N)$, where $W_1,\ldots,W_N$ are independent.
Suppose that $\E X=0$ and that the $j$th-coordinate sensitivity of $f$ is
at most $c_j$. Then, for every $p\ge2$,
\[
\norm{X}_{L_p}
\le
\left(
\frac p2\sum_{j=1}^N c_j^2
\right)^{1/2}.
\]
\end{lemma}

\begin{proposition}\label{prop:cube}
Let $h_i:\{-1,1\}^m\to\mathbb R$, $i\in[m]$, satisfy \(\E_i h_i=0\) and \( \E[h_i\mid\varepsilon_i]=0 \). Assume that for every $i\neq j$,
\[
\sup_{x\in\{-1,1\}^m} \abs{h_i(x)-h_i(x^{(j)})} \le\gamma.
\]
Then, for every $m\ge1$ and $p\ge2$, there is a universal constant $C$ such that
\[
\norm{\sum_{i=1}^m h_i}_{L_p} \le Cmp\gamma.
\]
\end{proposition}

\begin{proposition}\label{prop:transfer}
Let $h_1,\ldots,h_n$ be functions of arbitrary independent coordinates $Z_1,\ldots,Z_n$ satisfying
\begin{equation}\label{eq:double-center}
\E[h_i\mid Z_i]=0,
\qquad
\E[h_i\mid Z_{-i}]=0.
\end{equation}
Assume also that changing coordinate $Z_j$ changes $h_i$ by at most $\beta$ for every $j\neq i$.
Then, for every $p\ge2$,
\[
\norm{\sum_{i=1}^n h_i(Z)}_{L_p} \le 16pn\beta.
\]
\end{proposition}

\begin{proof}
Let $Z'=(Z'_1,\ldots,Z'_n)$ be an independent copy of $Z$ and set
\(
W_i=(Z_i^+,Z_i^-):=(Z_i,Z'_i), \ W=(W_1,\ldots, W_n).\)
We also introduce independent unbiased signs $\varepsilon_1,\ldots,\varepsilon_n$, independent of $W$, and use them to select one element from each pair $W_i$ by defining
\[
Z_i^\varepsilon=
\begin{cases}
Z_i^+,&\varepsilon_i=1,\\
Z_i^-,&\varepsilon_i=-1.
\end{cases}
\]
For each $i$, let $\varepsilon^{(i)}$ denote the sign vector obtained from $\varepsilon$ by flipping its $i$th coordinate. Conditionally on $W$, we then define
\[
Q_i(\varepsilon;W)
=
\frac12\left[
h_i(Z^\varepsilon)-h_i(Z^{\varepsilon^{(i)}})
\right].
\]

We first relate the original sum $\sum_i h_i(Z)$ to the randomized quantities $Q_i$. By the second centering identity in \eqref{eq:double-center},
\[ h_i(Z) = \E_{Z'_i} \left[ h_i(Z)-h_i(Z'_i,Z_{-i}) \mid Z \right]. \]
Therefore, Jensen's inequality yields
\begin{align}
\norm{\sum_{i=1}^nh_i(Z)}_{L_p} & \le \norm{ \sum_{i=1}^n \left[h_i(Z)-h_i(Z'_i,Z_{-i})\right]}_{L_p} \notag\\
& = 2\norm{\sum_{i=1}^nQ_i(\mathbf 1;W)}_{L_p}.
\label{eq:symmetrize}
\end{align}
The next observation allows us to replace the fixed orientation $\mathbf 1$ by random Rademacher signs. Fix a deterministic $\xi\in\{-1,1\}^n$ and swap the two entries of $W_j$ whenever $\xi_j=-1$. Since each pair $(Z_j,Z'_j)$ is exchangeable, this operation leaves the law of $W$ unchanged, while it transforms $Q_i(\mathbf1;W)$ into $Q_i(\xi;W)$. Averaging this identity over an independent uniform choice $\xi=\varepsilon$ therefore gives
\begin{equation}\label{eq:orientation}
\norm{\sum_iQ_i(\mathbf1;W)}_{L_p} = \norm{\sum_iQ_i(\varepsilon;W)}_{L_p}.
\end{equation}

We now examine the structure of $Q_i$ as a function of the Rademacher signs. For fixed $W$, flipping $\varepsilon_i$ reverses the sign of $Q_i$, so that $Q_i(\varepsilon^{(i)};W)=-Q_i(\varepsilon;W)$. Consequently, there exists a function $r_i(\varepsilon_{-i};W)$ such that
\[ Q_i(\varepsilon;W) = \varepsilon_i r_i(\varepsilon_{-i};W). \]
To impose the second centering condition required by Proposition~\ref{prop:cube}, we define
\[
b_i(W) = \E_{\varepsilon_{-i}} r_i(\varepsilon_{-i};W), \qquad \widetilde Q_i(\varepsilon;W) = Q_i(\varepsilon;W)-\varepsilon_i b_i(W).
\]
Conditionally on $W$, this definition ensures that
\[ \E[\widetilde Q_i\mid\varepsilon_{-i},W]=0, \qquad \E[\widetilde Q_i\mid\varepsilon_i,W]=0. \]
Moreover, for every $j\neq i$, the coordinate sensitivity assumption on $h_i$ implies
\[ \abs{Q_i(\varepsilon;W)-Q_i(\varepsilon^{(j)};W)} \le\beta.\]
Since the correction term $\varepsilon_i b_i(W)$ does not depend on $\varepsilon_j$ when $j\neq i$, subtracting it does not change this sensitivity bound. Proposition~\ref{prop:cube}, applied conditionally on $W$, therefore gives
\begin{equation}\label{eq:Qtilde}
\norm{\sum_{i=1}^n\widetilde Q_i}_{L_p(\varepsilon\mid W)} \le Cpn\beta.
\end{equation}

It remains to control the centering defect introduced by the terms $b_i(W)$. To make their dependence on $W$ explicit, fix a pair $W_i=(u,v)$ and define
\[
d_i(x_{-i};u,v)
=
\frac12
\left[
h_i(u,x_{-i})-h_i(v,x_{-i})
\right].
\]
With this notation, the definition of $b_i(W)$ can be rewritten as
\begin{equation}\label{eq:b-representation}
b_i(W)
=
\E_{\varepsilon_{-i}}
d_i(Z_{-i}^{\varepsilon};Z_i^+,Z_i^-).
\end{equation}
When we average over $W_{-i}$ and $\varepsilon_{-i}$, the random vector $Z_{-i}^{\varepsilon}$ has the same law as an independent draw from the original distribution of $Z_{-i}$. The first centering identity in \eqref{eq:double-center} therefore implies that $\E[b_i(W)\mid W_i]=0$. Moreover, if $j\neq i$, replacing the pair $W_j$ by another pair changes the right-hand side of \eqref{eq:b-representation} by at most $\beta$. Hence, applying Lemma~\ref{lem:bd} conditionally on $W_i$ gives
\[
\norm{b_i}_{L_p} \le \beta\sqrt{\frac{(n-1)p}{2}}.
\]

To aggregate these bounds over $i$, let \( B(W)=\left(\sum_i b_i(W)^2\right)^{1/2} \). Applying the triangle inequality in $L_{p/2}$ to $\sum_ib_i^2$ and then taking square roots yields
\[
\norm{B}_{L_p} \le \left(\sum_{i=1}^n\norm{b_i}_{L_p}^2\right)^{1/2}.
\]
Consequently,
\begin{equation}\label{eq:B-moment}
\norm{B}_{L_p} \le \beta\sqrt{\frac{n(n-1)p}{2}} \le n\beta\sqrt{\frac p2}.
\end{equation}

We next use this estimate to control the Rademacher sum generated by the centering defect. Conditionally on $W$, flipping the sign $\varepsilon_i$ changes $\sum_i\varepsilon_i b_i(W)$ by $2\abs{b_i(W)}$. A second application of Lemma~\ref{lem:bd} therefore gives
\[
\norm{\sum_i\varepsilon_i b_i(W)}_{L_p(\varepsilon\mid W)}
\le
\sqrt{2p}\,B(W).
\]
Combining this conditional estimate with \eqref{eq:B-moment} yields
\begin{equation}\label{eq:defect}
\norm{\sum_i\varepsilon_i b_i(W)}_{L_p} \le pn\beta.
\end{equation}
Hence, the decomposition of $Q_i$ into $\widetilde Q_i$ and the centering defect, together with \eqref{eq:Qtilde} and \eqref{eq:defect}, gives
\[
\norm{\sum_iQ_i(\varepsilon;W)}_{L_p} \le (C+1)pn\beta.
\]
Finally, Proposition~\ref{prop:cube}, with the choice
$c=1/(2e+2)$, gives $C<7$. Therefore, $2(C+1) < 16.$ Combining this estimate with \eqref{eq:symmetrize} and
\eqref{eq:orientation}, we obtain
\[
\norm{\sum_{i=1}^n h_i(Z)}_{L_p} \le 16pn\beta.
\]
This completes the proof.
\end{proof}

\end{document}
